% !PRINTED 174
\weberpage{174}

%<<CONT
and thus, by (16) of the preceding section, they yield
\[
\alpha s_0 + \beta s_1 + \gamma s_2 = 0,
\]
\[
\alpha^2 s_0 + 2\alpha\beta s_1 + 2\alpha\gamma s_2 + \beta^2 s_2 + 2\beta\gamma s_3 + \gamma^2 s_4 = 0,
\]
in which $s_0 = 3$. If one multiplies the latter equation by $s_0$ and subtracts from it the square of the first, it follows that
\weq{6}{
\beta^2(s_0 s_2 - s_1^2) + 2\beta\gamma(s_0 s_3 - s_1 s_2) + \gamma^2(s_0 s_4 - s_2^2) = 0.
}
The discriminant of this quadratic equation
\[
(s_0 s_3 - s_1 s_2)^2 - (s_0 s_2 - s_1^2)(s_0 s_4 - s_2^2)
\]
gives, when expanded,
\[
-s_0(s_0 s_2 s_4 - s_0 s_3^2 - s_1^2 s_4 - s_2^3 + 2 s_1 s_2 s_3),
\]
and is therefore, if $D$ denotes the discriminant of the cubic equation (1) [\S.\,46, (9)], equal to $-3D$; and the solution of (6) accordingly gives
\weq{7}{
\frac{\beta}{\gamma} = \frac{-(s_0 s_3 - s_1 s_2) + \sqrt{-3D}}{s_0 s_2 - s_1^2}.
}
With this, equation (3) has been reduced to a pure equation, which for its solution requires only the extraction of a cube root. Which sign we wish to give the square root occurring in (7) is left to our discretion.

To make use of the \gesperrt{Tschirnhausen} transformation for the biquadratic equation, one may proceed, for instance, by determining the $\alpha_0, \alpha_1, \alpha_2, \alpha_3$ in the transformed equation
\weq{8}{
y^4 + b_1 y^3 + b_2 y^2 + b_3 y + b_4 = 0
}
so that
\weq{9}{
b_1 = 0, \quad b_3 = 0
}
whereby (8) passes into a quadratic equation for $y^2$, from whose roots the square roots must still be extracted, so that, after the equations (9) have been solved, a square root must still be extracted twice more.

If, by means of the first equation (9), one eliminates $\alpha_0$ from the second, there arises a homogeneous cubic equation among $\alpha_1, \alpha_2, \alpha_3$. One may therefore assume one of the two ratios $\alpha_1:\alpha_2:\alpha_3$ arbitrarily, e.g. set $\alpha_3 = 0$, and obtains for the determination of the other a \gesperrt{cubic equation}. By this, then, the solution of the biquadratic equation is reduced to that of a cubic and to square roots.
